% Part: incompleteness
% Chapter: introduction
% Section: definitions

\documentclass[../../../include/open-logic-section]{subfiles}

\begin{document}

\olfileid{inc}{int}{def}

\olsection{تعریفاں}

ہلبرٹ دے منصوبے وچ ریاضی نوں رسمی شکل دین،
تے فیر ایس رسمی نظام دی بے تضادی تے کاملیت
دکھاون لئی، پہلا کم اک زبان، منطقی ڈھانچہ تے
مسلّماں دا نظام چُننا ہووے گا۔ ایتھے اسی فرض
کر لئیے کہ ریاضی نوں پہلے درجے دی زبان وچ
رسمی شکل دتی جا سکدی اے۔ مطلب ایہ کہ
کوئی ایہو جیہا سیٹ اے جیہدے وچ !!{constant}s،
!!{function}s تے !!{predicate}s ہون۔ ایہ نشان
پہلے درجے دی منطق دے رابطیاں تے مسوّراں نال
مل کے ریاضی دے دعوے بیان کر سکدے نیں۔ اکثر لوک
مندے نیں کہ ایڈی زبان موجود اے: سیٹ نظریے
دی زبان، جیہدے وچ $\in$ واحد غیر منطقی نشان
اے۔ پہلی نظر وچ ایڈی سادی زبان دے اِنّی
وڈی اظہار دی طاقت رکھن دا دعویٰ ناقابلِ
یقین لگدا اے؛ ایہ ثابت کرن لئی بہت کم
کرنا پیا کہ عملاً ساری ریاضی ایس بہت
محدود لفظیات وچ بیان ہو سکدی اے۔ گل
سادی رکھن لئی فی الحال بحث نوں حساب
تک محدود کریے، یعنی ریاضی دا اوہ
حصہ جیہڑا صرف قدرتی عدداں~$\Nat$ نال
سروکار رکھدا اے۔ حساب دے حقائق بیان
کرن دی فطری زبان~$\Lang L_A$ اے۔
$\Lang L_A$ وچ اک دو تھانواں
والا !!{predicate}~$<$، اک
!!{constant}~$\Obj 0$،
اک یک-مقامی !!{function}~$\prime$،
تے دو دو تھانواں والے
!!{function}s~$+$ تے~$\times$ نیں۔

\begin{defn}
!!{sentence}s نال بنیا سیٹ~$\Gamma$ اک
\emph{نظریہ} اے جے اوہ منطقی نتیجیاں
دے لحاظ نال بند ہووے، یعنی جے
$\Gamma = \Setabs{!A}{\Gamma \Entails
!A}$۔
\end{defn}

نظریہ متعین کرن دے دو سوکھے طریقے
نیں۔ اک ایہ کہ جیہڑے !!{sentence}s کسے
!!{structure} وچ سچ ہون، اوہناں دا سیٹ لے لئیے۔
مثال لئی $\Lang L_A$ دی
اوہ !!{structure} ویکھو جیہدے
!!{domain}~$\Nat$ اے تے
ہر غیر منطقی نشان دی تعبیر
اوہو اے جیہڑی توقع ہوندی اے۔

\begin{defn}\ollabel{def:standard-model}
\emph{حساب دا معیاری نمونہ}
اوہ !!{structure}~$\Struct
N$ اے جیہدی تعریف ایویں اے:
\begin{enumerate}
\item $\Domain N = \Nat$
\item $\Assign{\Obj 0}{N} = 0$
\item $\Assign{\Obj \prime}{N}(n) = n + 1$ ہر $n \in \Nat$ لئی
\item $\Assign{\Obj +}{N}(n, m) = n + m$ ہر $n, m \in \Nat$ لئی
\item $\Assign{\Obj \times}{N}(n, m) = n\cdot m$ ہر $n, m \in \Nat$ لئی
\item $\Assign{\Obj <}{N} = \Setabs{\tuple{n, m}}{n \in \Nat, m \in
  \Nat, n < m}$
\end{enumerate}
\end{defn}

'$\times$' تے '$\cdot$' دا فرق دھیان
نال ویکھو: $\times$ حساب دی زبان
دا اک نشان اے۔ بے شک اسی ایہ نشان
ضرب یاد کراؤن لئی چُنیا اے، پر
$\times$ خود ضرب دا عمل نہیں؛ ایہ
دو تھانواں والا فنکشن نشان اے
(باقاعدہ ناں~$\Obj f^2_1$)۔
ایہدے برعکس $\cdot$ عام ضرب
دا \emph{فنکشن اے}۔ جدوں
$n \cdot m$ ویکھو تاں مراد
عدداں~$n$ تے~$m$ دا حاصلِ
ضرب اے؛ جدوں $x \times y$
ویکھو تاں گل حساب دی زبان
دے اک ٹرم بارے اے۔ معیاری
نمونے وچ فنکشن نشان~$\times$
دی تعبیر قدرتی عدداں اُتے
فنکشن~$\cdot$ نال کیتی
جاندی اے۔ جمع لئی اسی $+$
نوں حساب دی زبان دے فنکشن
نشان تے قدرتی عدداں اُتے
جمع دے فنکشن، دوناں لئی
ورتدے آں۔ ایتھے سیاق
توں طے کرنا پیندا اے کہ
مراد کیہ اے۔

\begin{defn}
\emph{سچّے حساب} دا نظریہ
اوہ سیٹ اے جیہدے وچ اوہی !!{sentence}s
نیں جیہڑے حساب دے معیاری
نمونے وچ پورے اتردے نیں،
یعنی
\[
\Th{TA} = \Setabs{!A}{\Sat{N}{!A}}.
\]
\end{defn}

$\Th{TA}$ اک نظریہ اے:
جے $\Th{TA} \Entails !A$
ہووے تاں $!A$ ہر اُس
!!{structure} وچ پورا
اتردا اے جیہڑی~$\Th{TA}$
نوں پورا کردی اے۔
$\Sat{N}{\Th{TA}}$ ہون
کارن~$\Sat{N}{!A}$
وی اے، تے ایس لئی
$!A \in \Th{TA}$۔

نظریہ~$\Gamma$ متعین
کرن دا دوجا طریقہ ایہ
اے کہ !!{sentence}s نال
بنے کسے سیٹ~$\Gamma_0$
دے سارے منطقی نتیجیاں
دا سیٹ لے لئیے۔ ایس
صورت وچ $\Gamma$،
$\Gamma_0$ دی منطقی
نتیجیاں دے لحاظ نال
«بندش» اے۔ ایہ طریقہ
تاں ہی کارآمد اے جے
$\Gamma_0$ صاف طور اُتے
متعین ہووے، مثلاً
اوہ~$\Gamma_0$ دے !!{element}s دی
فہرست دتی گئی ہووے۔
گھٹ توں گھٹ~$\Gamma_0$
دے بارے فیصلہ ہو سکنا
چاہیدا اے: کسے
!!a{sentence} دے
$\Gamma_0$ دا رکن
ہون یا نہ ہون دی
قابلِ حساب جانچ ہووے۔
$\Gamma_0$ دے جیہڑے
!!{sentence}s نیں، اوہناں نوں
$\Gamma$ دے \emph{مسلّمے}
آکھدے آں، تے آکھدے
آں کہ $\Gamma$،
$\Gamma_0$ نال
\emph{مسلّماتی طور متعین}
اے۔

\begin{defn}
نظریہ~$\Gamma$،
$\Gamma_0$ نال
\emph{مسلّماتی طور متعین}
اے جے تے صرف جے
\[
\Gamma = \Setabs{!A}{\Gamma_0 \Entails !A}
\]
\end{defn}

\begin{defn}
اگلے جملیاں نال مسلّماتی
طور متعین نظریے~$\Th{Q}$
نوں «رابنسن دا~$\Th{Q}$»
آکھدے نیں؛ ایہ حساب
دا بہت سادہ نظریہ اے۔
\begin{align*}
& \lforall[x][\lforall[y][(\eq[x'][y'] \lif \eq[x][y])]] \tag{$!Q_1$}\\
& \lforall[x][\eq/[\Obj 0][x']] \tag{$!Q_2$}\\
& \lforall[x][(\eq[x][\Obj 0] \lor \lexists[y][\eq[x][y']])] \tag{$!Q_3$}\\
& \lforall[x][\eq[(x + \Obj 0)][x]] \tag{$!Q_4$}\\
& \lforall[x][\lforall[y][\eq[(x + y')][(x + y)']]] \tag{$!Q_5$}\\
& \lforall[x][\eq[(x \times \Obj 0)][\Obj 0]] \tag{$!Q_6$}\\
& \lforall[x][\lforall[y][\eq[(x \times y')][((x \times y) + x)]]] \tag{$!Q_7$}\\
& \lforall[x][\lforall[y][(x < y \liff \lexists[z][\eq[(z' + x)][y]])]] \tag{$!Q_8$}
\end{align*}
!!{sentence}s نال بنیا سیٹ
$\{!Q_1, \dots, !Q_8\}$
ہی $\Th{Q}$ دے مسلّمے
نیں؛ ایس لئی $\Th{Q}$
اوہناں توں نکلن والے
سارے !!{sentence}s
دا سیٹ اے:
\[
\Th{Q} = \Setabs{!A}{\{!Q_1, \dots, !Q_8\} \Entails !A}.
\]
\end{defn}

\begin{defn}
فرض کرو $!A(x)$،
$\Lang L_A$ دا
!!a{formula} اے جیہدے
آزاد متغیر~$x$ تے
$y_1$, \dots, $y_n$
نیں۔ فیر ایس شکل دا
ہر !!{sentence}
\[
\lforall[y_1][\dots\lforall[y_n][((!A(\Obj 0) \land \lforall[x][(!A(x)
\lif !A(x'))]) \lif \lforall[x][!A(x)])]]
\]
\emph{استقرائی خاکے}
دی اک مثال اے۔

\emph{پیانو حساب}~$\Th{PA}$
اوہ نظریہ اے جیہڑا
$\Th{Q}$ دے مسلّماں
تے استقرائی خاکے
دیاں ساریاں مثالاں
نال مسلّماتی طور
متعین ہوندا اے۔
\end{defn}

\begin{explain}
استقرائی خاکے دی ہر
مثال~$\Struct{N}$
وچ سچ اے۔ ایہ گل
سب توں سوکھی اوہدوں
ویکھی جا سکدی اے
جدوں !!{formula}~$!A$
دا صرف اک آزاد
!!{variable}~$x$
ہووے۔ فیر $!A(x)$
اک ذیلی سیٹ~$X_{!A}$
متعین کردا اے؛
ایہ $\Nat$ دا
ذیلی سیٹ اے،
تے تعریف
$\Struct{N}$ وچ اے۔
$X_{!A}$ وچ
اوہ سارے~$n \in \Nat$
نیں جنہاں لئی
$\Sat{N}{!A(x)}[s]$
ہووے، جدوں
$s(x) = n$۔
استقرائی
خاکے دی متعلقہ
مثال ایہ اے:
\[
((!A(\Obj 0) \land \lforall[x][(!A(x) \lif !A(x'))]) \lif 
  \lforall[x][!A(x)]).
\]
جے ایس دا مقدمہ
$\Struct{N}$ وچ
سچ ہووے تاں
$0 \in X_{!A}$
تے جدوں وی
$n \in X_{!A}$
ہووے، $n+1$
وی اوہدے وچ
ہووے گا۔
$0 \in X_{!A}$
توں~$1 \in X_{!A}$
ملدا اے؛
$1 \in X_{!A}$
توں~$2 \in X_{!A}$
ملدا اے؛ تے
ایویں اگے۔
ایس لئی ہر
$n \in \Nat$
لئی~$n \in X_{!A}$
اے۔ ایہو
گل اے کہ
$\lforall[x][!A(x)]$
، $\Struct{N}$
وچ پورا اتردا اے۔
\end{explain}

$\Th{Q}$ تے $\Th{PA}$
دونیں مسلّماتی طور
متعین نظریے نیں۔
وڈا سوال ایہ اے
کہ اوہ کِنّے طاقتور
نیں؟ مثال لئی،
کیہ $\Th{PA}$،
$\Nat$ بارے
$\Lang L_A$ وچ
بیان ہو سکّن
والیاں ساریاں
سچائیاں ثابت کر
سکدا اے؟ ٹھیک
ٹھیک پچھیے تاں
کیہ $\Th{PA}$
دے مسلّمے
$\Lang L_A$ وچ
پوچھے جا سکّن
والے ہر سوال
دا فیصلہ کردے
نیں؟

ایہی سوال ہور
طرح ایویں پچھ
سکدے آں: کیہ
$\Th{PA} = \Th{TA}$
اے؟ ظاہر اے کہ
$\Th{TA}$،
$\Nat$ بارے
ساریاں سچائیاں
ثابت کردا اے
(یعنی اوہناں
نوں شامل کردا
اے) تے
$\Lang L_A$
وچ بیان ہو
سکّن والے
سارے سوالاں
دا فیصلہ کردا
اے۔ وجہ ایہ
اے کہ جے
$!A$،
$\Lang L_A$
دا !!a{sentence}
ہووے، تاں یا
$\Sat{N}{!A}$
یا
$\Sat{N}{\lnot !A}$
؛ ایس لئی
یا
$\Th{TA} \Entails !A$
یا
$\Th{TA} \Entails \lnot !A$۔
ایہو جیہے
نظریے نوں
\emph{!!{complete}}
آکھدے نیں۔

\begin{defn}
نظریہ~$\Gamma$
\emph{!!{complete}}
اے جے تے صرف
جے اوہدی زبان
دے ہر
!!{sentence}~$!A$
لئی یا
$\Gamma \Entails !A$
یا
$\Gamma \Entails \lnot !A$
ہووے۔
\end{defn}

\begin{explain}
تمامیت دے قضیے
موجب،
$\Gamma \Entails !A$
جے تے صرف جے
$\Gamma \Proves
!A$؛ ایس لئی
$\Gamma$ کامل
اے جے تے
صرف جے اوہدی
زبان دے ہر
!!{sentence}~$!A$
لئی یا
$\Gamma \Proves !A$
یا
$\Gamma \Proves \lnot !A$
ہووے۔
\end{explain}

اک ہور سوال وی
اگے آندا اے:
کیہ کوئی حسابی
طریقہ اے جیہدے
نال اسی پرکھ
سکیے کہ کوئی
!!a{sentence}،
$\Th{TA}$ وچ
اے، $\Th{PA}$
وچ اے، یا
صرف~$\Th{Q}$
وچ اے؟
کسے سیٹ
(مثلاً
!!{sentence}s نال بنے سیٹ)
دے !!{decidable}
ہون دی تعریف
دے کے ایس
سوال نوں ہور
ٹھیک بنا
سکدے آں۔

\begin{defn}
سیٹ~$X$
\emph{!!{decidable}}
اے جے تے
صرف جے اک
حسابی طریقہ
ہووے جیہڑا
انپٹ~$x$
اُتے~$1$
واپس دیوے
جے $x \in X$
ہووے، تے نہیں
تاں~$0$ دیوے۔
\end{defn}

ہن سوال ایہ
بندا اے:
کیہ $\Th{TA}$
($\Th{PA}$،
$\Th{Q}$)
!!{decidable}
اے؟

ایہناں سارے
سوالاں دا
جواب «نہیں»
ہووے گا۔
ایہناں وچوں
کوئی وی نظریہ
فیصلہ کرن جوگ
نہیں۔ پر ایہ
واقعہ صرف ایہناں
خاص نظریاں
تک محدود نہیں۔
حقیقت وچ،
کجھ شرطاں
پوری کرن والا
\emph{ہر}
نظریہ ایہناں
نتیجیاں ہیٹھ
آندا اے۔ اک
شرط، جیہڑی
$\Th{Q}$ تے
$\Th{PA}$ پوری
کردے نیں،
ایہ اے کہ
اوہ مسلّماں
دے اک سیٹ نال
متعین نیں،
جیہڑا !!a{decidable}
اے۔

\begin{defn}
نظریہ
\emph{!!{axiomatizable}} اے جے اوہ
مسلّماں دے اک سیٹ نال
مسلّماتی طور متعین ہووے،
تے اوہ سیٹ !!a{decidable} ہووے۔
\end{defn}

\begin{ex}
جیہڑا نظریہ
!!{sentence}s نال بنے
محدود سیٹ
نال مسلّماتی
طور متعین اے،
اوہ
!!{axiomatizable}
اے؛ وجہ ایہ
کہ ہر محدود
سیٹ
!!{decidable}
ہوندا اے۔
ایس لئی، مثال
دے طور اُتے،
$\Th{Q}$
!!{axiomatizable}
اے۔

$\Th{PA}$
ورگے خاکیاں
راہیں متعین
نظریے وی
!!{axiomatizable}
نیں۔ ایہ
پرکھّن لئی
کہ $!B$،
$\Th{PA}$
دے مسلّماں
وچ شامل اے
یا نہیں، یعنی
اوہ فنکشن
$\Char{X}$
حساب کرن لئی
جیہدے وچ
$\Char{X}(!B) = 1$
ہووے جے
$!B$،
$\Th{PA}$
دا مسلمہ
اے؛ نہیں
تاں~$= 0$
ہووے،
ایہ طریقہ
ورت سکدے آں۔
پہلاں پرکھو
کہ $!B$،
$\Th{Q}$
دے مسلّماں
وچوں اک اے
یا نہیں۔
جے اے،
تاں جواب
«ہاں» تے
$\Char{X}(!B) = 1$
اے۔ جے
نہیں، تاں
پرکھو کہ
ایہ استقرائی
خاکے دی
اک مثال
اے یا نہیں۔
ایہ کم
باقاعدہ طریقے
نال ہو
سکدا اے؛
ایس معاملے
وچ شاید
ایویں ویکھنا
سب توں سوکھا
اے: استقرائی
خاکے دی
ہر مثال
چند کلی
مسوّراں توں
شروع ہوندی
اے، جنہاں
مگروں اک
شرطیے والا
ذیلی
!!{formula}
آندا اے۔
اُس شرطیے
دا تالی
$\lforall[x][!A(x, y_1, \dots, y_n)]$
اے۔ ایتھے
$x$ تے
$y_1$, \dots, $y_n$
سارے~$!A$
دے آزاد
متغیر نیں،
تے $!B$
دے شروع والے
مسوّرے
$y_1$, \dots,~$y_n$
نوں باندھدے
نیں۔ ایہ~$!A$
کڈھن تے ایہدے
آزاد متغیراں
نوں شروع دے کلی
مسوّراں تے
$\lforall[x]$
نال بندھے
متغیراں
نال ملاؤن
مگروں، اسی
پرکھدے آں
کہ شرطیے
دا مقدمہ
ایس شکل
نال ملدا
اے یا نہیں:
\[
!A(\Obj 0, y_1, \dots, y_n) \land \lforall[x][(!A(x, y_1, \dots, y_n)
\lif !A(x', y_1, \dots, y_n))]
\]
جے ملدا اے
تاں $!B$
استقرائی
خاکے دی
مثال اے؛
جے نہیں ملدا
تاں $!B$
ایہو جیہی مثال
نہیں۔
\end{ex}

ایس سوال،
تے زیادہ
عمومی سوال
کہ کیہڑے
نظریے
کامل یا
فیصلہ کرن
جوگ نیں،
دے جواب
لئی اگلی
تعریف وی
فائدہ مند
ہووے گی۔
یاد کرو
کہ سیٹ~$X$
!!{enumerable}
اے جے تے
صرف جے اوہ
خالی اے
یا کوئی
!!a{surjective}
فنکشن~$f \colon \Nat
\to X$ موجود
اے۔ ایسے
فنکشن نوں
$X$ دی
گنتی وار
فہرست آکھدے
نیں۔

\begin{defn}
سیٹ~$X$
\emph{!!{computably enumerable}}
(مختصراً
!!{c.e.}) اے
جے تے
صرف جے اوہ
خالی ہووے
یا اوہدی
قابلِ حساب
گنتی وار
فہرست
موجود ہووے۔
\end{defn}

!!{axiomatizability}
دے نال
اک ہور
شرط جیہڑی
نامکملیت
دے قضیاں
لئی متعلقہ
نظریاں اُتے
لاگو ہووے
گی، ایہ
اے کہ اوہ
قابلِ حساب
فنکشناں
تے
!!{decidable}
تعلقّاں بارے
بنیادی حقائق
ثابت کرن
جوگے ہون۔
«بنیادی
حقائق»
توں مراد
اوہ
!!{sentence}s
نیں جیہڑے
قابلِ حساب
فنکشناں
دے
مدخلاں
لئی قدراں
بیان کردے
نیں۔
«کافی طاقتور»
توں مراد
ایہ اے
کہ متعلقہ
نظریے
ایہ جملے
اپنے قضیاں
وچ شامل
کردے ہون۔
مثال لئی
اک بنیادی
قابلِ حساب
فنکشن، جمع،
نوں ویکھو۔
$+$ دے
مدخلاں
$2$ تے
$3$ لئی
قدر~$5$
اے، یعنی
$2+3 = 5$۔
حساب دی
زبان وچ
$+$ دے
مدخلاں
$2$ تے
$3$ لئی
قدر~$5$
ہون دا
بیان ایہ
جملہ کردا اے:
$(\num{2} + \num{3}) = \num{5}$۔
مثلاً
$\Th{Q}$
ایہ جملہ
ثابت کردا
اے۔ عام
طور اُتے
اسی چاہندے
آں کہ
ہر قابلِ
حساب فنکشن
$f(x_1, x_2)$
لئی کوئی
!!a{formula}~$!A_f(x_1, x_2, y)$
زبان~$\Lang{L_A}$
وچ ایویں ہووے کہ
$\Th{Q}
\Proves !A_f(\num{n_1}, \num{n_2}, \num{m})$
جدوں وی
$f(n_1, n_2) =
m$ ہووے۔
ایس طرح
$\Th{Q}$
ثابت کردا
اے کہ
$f$ دی
قدر مدخلاں
$n_1$،
$n_2$
لئی~$m$
اے۔ حقیقت
وچ اسی
ایہدے توں
کجھ ودھ
وی منگدے
آں: نظریہ
ایہ وی
ثابت کرے
کہ~$f$ دی
قدر مدخلاں
$n_1$،
$n_2$ لئی
کوئی
ہور عدد
نہیں ہو
سکدی۔
!!{decidable}
تعلقّاں
لئی وی
ایویں
اے۔ اگلی
دو تعریفاں
ایہ گل
ٹھیک
بناندیاں
نیں۔

\begin{defn}
!!{formula}~$!A(x_1, \dots, x_k, y)$
فنکشن~$f\colon \Nat^k \to \Nat$
دی نظریے~$\Gamma$ وچ
\emph{!!{represents}} کردا اے
جے تے
صرف جے،
جدوں وی
$f(n_1,
\dots, n_k) = m$
ہووے، تاں
\begin{enumerate}
\item $\Gamma \Proves !A(\num{n_1}, \dots, \num{n_k}, \num{m})$، تے
\item $\Gamma \Proves \lforall[y](!A(\num{n_1}, \dots, \num{n_k},
y) \lif y = \num{m})$۔
\end{enumerate}
\end{defn}

\begin{defn}
% PNB-CMP-047: source relation definition omits the theory Gamma used in its two conditions.
!!{formula}~$!A(x_1, \dots, x_k)$
نظریے~$\Gamma$ وچ
تعلق~$R \subseteq \Nat^k$
دی \emph{!!{represents}} کردا اے
جے تے
صرف جے،
\begin{enumerate}
\item جدوں وی $R(n_1, \dots, n_k)$ ہووے،
  $\Gamma \Proves !A(\num{n_1},
\dots, \num{n_k})$، تے
\item جدوں وی $R(n_1, \dots, n_k)$ نہ ہووے،
  $\Gamma \Proves \lnot
!A(\num{n_1}, \dots, \num{n_k})$۔
\end{enumerate}
\end{defn}

نامکملیت دے
قضیے لاگو
ہون لئی
نظریہ «کافی
طاقتور» اے
جے اوہ
سارے قابلِ
حساب فنکشناں
تے سارے
!!{decidable}
تعلقّاں دی
!!{represents}
کردا اے۔
$\Th{Q}$~
تے اوہدے
توسیعی نظریے
ایہ شرط پوری
کردے نیں،
پر ایس
گل دا ثبوت
دین لئی
سانوں
ہن کجھ ویلا
لگے گا۔
ایہ
$\Th{Q}$
دے ثبوتی
اختیار بارے
غیر معمولی
حقیقت اے؛
ایہ ثابت
کرنی اوکھی
اے کیوں
جے
$\Th{Q}$
دے مسلّمے
گھٹ نیں
تے سارے
حقائق اوہناں
توں اخذ
کرنے پینے
نیں۔ پر
$\Th{Q}$
خود بہت
کمزور نظریہ
اے۔ ایس لئی
بھاویں ایہ
ثابت کرنا
اوکھا اے
کہ
$\Th{Q}$
سارے قابلِ
حساب فنکشناں
دی نمائندگی
کردا اے،
اکثر دلچسپ
نظریے
$\Th{Q}$
توں طاقتور
نیں، یعنی
اوہ
$\Th{Q}$
توں ودھ
کجھ ثابت
کردے نیں۔
جے
$\Th{Q}$
کجھ ثابت
کردا اے
تاں ہر
زیادہ طاقتور
نظریہ وی
اوہ ثابت
کردا اے۔
ایس لئی
$\Th{Q}$
سارے قابلِ
حساب فنکشناں
دی نمائندگی
کردا ہووے
تاں ہر
زیادہ طاقتور
نظریہ وی
ایہ کردا
اے۔ ایہ
دکھاون نال
نامکملیت
دے قضیے
بہت سارے
دلچسپ نظریاں
اُتے لاگو
ہوندے نیں۔
پوری ریاضی
نوں رسمی
شکل دین
دا امیدوار
نظریہ یقیناً
$\Th{Q}$
دا ثابت
کیتا سب کجھ
ثابت کرے؛
گھٹ توں
گھٹ اوہنوں
ابتدائی
حساباں دے
نتیجے سنبھال
سکنے چاہیدے
نیں۔ سو
«ساری ریاضی»
دے نظریے
دا کوئی وی
امیدوار اوہناں
وچوں ہووے
گا جنہاں
اُتے نامکملیت
دے قضیے
لاگو ہوندے
نیں۔

\end{document}
